\FloatBarrier
\chapter{Conclusions and Future Outlook}
\label{chap:conclusions}

This thesis has provided a comprehensive overview of the theoretical foundations
and practical implementations of atomic and quantum sensors, with a specific
emphasis on the development of laser-written vapor cells. Central to this work
is the exploration of the FLICE (Femtosecond Laser Irradiation followed by
Chemical Etching) technique. This cutting-edge manufacturing method enables the
fabrication of complex, arbitrary three-dimensional cavities within fused silica,
overcoming the geometric limitations inherent to conventional glass-blowing and
machining methods.

To ensure the utmost purity of the internal sensing gas ---a critical requirement
for preserving atomic coherence lifetimes--- adhesive-free sealing techniques were
extensively analyzed, particularly optical contact bonding. By eliminating
outgassing contaminants at the interface, this bonding approach facilitates
exceptional sensor resolutions that closely approach the fundamental quantum
limit. The successful implementation of these pristine vapor cells is crucial
for achieving sensitivities on the order of attoteslas for Optically Pumped
Magnetometers (OPMs), picovolts (pV) for Rydberg electrometers, and femtohertz
for advanced atomic clocks.

Furthermore, a rigorous manufacturing workflow demands precise surface
preparation and meticulous validation. Through an extensive literature review,
this work synthesized various methods for the critical cleaning and surface
activation of fused silica substrates. The analyzed protocols encompassed both
traditional wet-chemical approaches and, more recently, plasma-based activation
techniques. To complete the fabrication pipeline, a suite of characterization
methods was detailed to monitor each stage of the process. This included
surface integrity evaluations utilizing Newton rings and Atomic Force Microscopy
(AFM), surface energy assessments via contact angle goniometry, and mechanical
robustness evaluations through razor blade insertion and tensile testing.

The theoretical frameworks and fabrication methodologies detailed in the early
chapters were ultimately translated into practical applications during an
internship at QSENSATO Srl. These hands-on validation stages included the
experimental execution of optical contact bonding trials to test seal viability,
the design of a magnetic corrosion inspection system for railway tracks, and
the successful implementation of a Zero-Field Resonance experiment.

Ultimately, the synergy of advanced laser micromachining and robust,
adhesive-free bonding paves the way for a new generation of miniaturized, highly
scalable quantum sensors. The progress and methodologies demonstrated throughout
this thesis underscore the immense potential of this technology to transition
from highly controlled laboratory environments into robust, real-world field
applications. As manufacturing techniques continue to mature and become more
accessible, these quantum devices will undoubtedly drive unprecedented
advancements in precision metrology, medical diagnostics, and widespread
industrial monitoring, opening a bright and transformative future for quantum
sensing technologies.

I would
like to express my deepest gratitude to my supervisors and the entire team at
QSENSATO Srl, whose invaluable guidance, support, and collaborative expertise
were fundamental to the successful completion of this thesis.